% Insertion-only repair; no documentclass and no new cross-reference labels.
% Target main SHA256:
% ebb86167e4a292d01058ac8a991d9ad4b804d403cc1787ced080a3c6a73350f8
% In imbalanced-chen-transfer.tex replace lines 2002--2027, beginning
% "Choose a single nonnegative smooth outer approximation of each"
% and ending "in a shrinking boundary strip."
% The same source passage is uniform-prime-products-upper.tex:99--124.
% In particular replace the unjustified blanket sentence at main 2019--2020:
% "prime-product integral has bounded density; the finite monomial
% and ordering boundaries have measure zero."
% Do not replace the following exact-coprime distribution paragraph.

\paragraph{Prime-product mass lemma and radial boundaries.}
The mass calculation requires both exponent coordinates and the
radial coordinate $y=m/N$. For a fixed number $k$ of prime slots put
\[
 t_j=\frac{\log p_j}{\log N}\quad(j<k),\qquad
 t_k^*=1-\sum_{j<k}t_j,\qquad y=\frac{p_1\cdots p_k}{N}.
\]
Fix $0<a_0<a$, and work on an ambient domain
$t_j,t_k^*\ge a_0$, $0<y\le C$, with $C$ fixed.
The normalized labelled prime-product mass converges, on continuous
test weights and subsequently on domains with null boundary in these
coordinates, to
\[
 d\mu_k=\frac{dt_1\cdots dt_{k-1}\,dy}
 {t_1\cdots t_{k-1}t_k^*}.
\]
For $k=1$ this means $d\mu_1=dy$. Fixed bounded constants attached
to the finite labels are retained. This is a bounded density on the
displayed simplex times a radial interval; it is not a density in
the full space of the $k$ prime exponents.

Here is the requisite proof and uniform error estimate. First take
$\eta\le y\le C$ with $\eta>0$ fixed. With the first $k-1$ primes
held fixed and $P=p_1\cdots p_{k-1}$, write $p_k=Ny/P$. The prime
number theorem and partial summation give the last-prime measure
\[
 \frac{dp_k}{\log p_k}
 =\frac{N}{P\log N}\frac{dy}{t_k^*+\log y/\log N},
\]
with error $O_A((N/P)(\log N)^{-A})$ for fixed smooth seminorms.
All nonempty relevant last-prime intervals have their endpoints at
least a fixed positive power of $N$. Since
\[
 \sum_{N^{a_0}/2\le p\le CN}\frac1p=O_{a_0,C}(1),
\]
the sum of these errors over the other primes is
$O_{A,k,a_0}(N(\log N)^{-A})$.
The reciprocal-prime measures in the remaining coordinates converge
to $dt_j/t_j$, again by prime-number-theorem partial summation.
On $y\ge\eta$ the displayed denominator tends uniformly to $t_k^*$,
proving the assertion first for smooth and then continuous weights.
The elementary upper bound for primes also gives
\[
 \#\{(p_1,\ldots,p_k):p_j\ge N^{a_0}/2,\ 
                  p_1\cdots p_k\le\eta N\}
 \ll_{k,a_0}\frac{\eta N}{\log N}.
\]
Indeed fix $P$, bound its last-prime count by
$O_{a_0}(\eta N/(P\log N))$, and sum reciprocal products.
This removes the radial restriction by first taking $N\to\infty$
and then $\eta\downarrow0$.

For completeness, each fixed monomial constraint
$\prod_{j=1}^kp_j^{\beta_j}\le cN^\gamma$ becomes exactly
\[
 \sum_{j<k}(\beta_j-\beta_k)t_j+\beta_k-\gamma
 \le\frac{\log c-\beta_k\log y}{\log N}.
\]
If the $\beta_j$ are not all equal, its limiting boundary is a
nonconstant affine hyperplane in the simplex coordinates and has
$\mu_k$-measure zero. If all $\beta_j=\lambda$ but
$\gamma\ne\lambda$, its truth value is eventually constant on
$\eta\le y\le C$. If all $\beta_j=\gamma=\lambda\ne0$, it is
the exact radial condition $y^\lambda\le c$ and must be smoothed
in $y$. When $\lambda=0$ its constant truth value $1\le c$ is
kept exactly. This classification covers ordering constraints,
equalities, and finite intersections and unions. In particular it
does not discard a pure total-product boundary by using the full
exponent-space measure. For example the empty one-prime domain
$p=N$ at even $N$ has exponent outer strips
$N^{1-\delta}\le p\le N$ of normalized limiting mass one, whereas
radial strips around $p/N=1$ have mass tending to zero with their width.

\paragraph{A positive outer approximation before sieving.}
For each fixed error budget $\tau>0$, choose $\eta$ sufficiently
small. On $\eta\le y\le C$ smooth the nonconstant affine walls
outward by small fixed exponent widths, and the critical pure-product
walls by small fixed radial widths. Retain a nonnegative smooth
majorant of the whole strip $0<y\le\eta$, supported in $y<2\eta$;
its mass is $O(\eta)N/\log N$. The preceding classification and
mass lemma then give one nonnegative bounded global outer weight
$W_{N,\tau}$ with
\[
 W_{N,\tau}\ge\mathbf1_{\mathcal D_N},\qquad
 \limsup_{N\to\infty}\frac{\log N}{N}
       \sum_{p_1,\ldots,p_k}
        (W_{N,\tau}-\mathbf1_{\mathcal D_N})\le\tau.
\]
Keep a separate lower-cutoff factor equal to one for $p_j\ge N^a$
and zero for $p_j\le N^a/2$. Here this factor-two transition has
vanishing excess mass: for $k\ge2$ its limiting face is a
nonconstant simplex wall, and for $k=1$ the cutoff $a<1$ is
noncritical. For a general critical pure-product cutoff the radial
rule is required instead of a fixed factor-two transition.
Thus every prime in the actual outer family remains at least $N^a/2$.

The operators $p_j\partial_{p_j}$ send $t_i$ to
$\delta_{ij}/\log N$ and send $y$ to $y$. Hence all normalized
derivatives on dyadic charts are bounded for each fixed smoothing
precision, uniformly in $N$ and the chart location. A smooth dyadic
partition of unity gives $O((\log N)^k)$ rectangles and sums back
to this one global majorant. The excess mass is therefore counted
once; only distribution errors acquire the harmless logarithmic
number of rectangles, absorbed by a larger requested saving.

The order of use is essential. First dominate the desired positive
prime count by its count with $W_{N,\tau}$. Then apply the smooth
exact-coprime distribution statement and the upper sieve below.
Small excess total mass alone is not asserted to control a difference
of signed progression remainders. After these estimates, take
$N\to\infty$ for the fixed family, widths, and factorization margins,
and only then take $\tau\downarrow0$. No averaging over $N$ or
uniformity in a cutoff tending to zero is used.
